\chapter{Проектирование автоматов и очереди FIFO}
\label{ch:fsmdesign}

В главе~\ref{ch:sync} мы познакомились с конечными автоматами и разобрали
пример детектора. Теперь научимся проектировать их систематически: от
словесного описания до схемы, и убедимся, что этот же рецепт работает для
любой управляющей части устройства. Во второй половине главы построим на тех
же приёмах очередь FIFO, которая нужна в любой системе, где скорости источника
и приёмника данных различаются.

Вопрос главы: как не угадывать автомат, а выводить его из словесного условия,
чтобы не забыть ни одного случая? Заодно мы соберём главный этап сквозного
проекта: автомат режимов «Охранной панели».

\begin{buildbox}
Вкладка \emph{Обучение}, книга «из блоков», глава «Автоматы», урок «Автомат на
двух состояниях» (затем «Каркас автомата»). Соберите автомат с состояниями
«выкл» и «вкл», переключаемый нажатием. Во вкладке \emph{Диаграммы} убедитесь,
что состояние меняется только по фронту такта. Запишите на листе, как вы
определили, \emph{сколько} состояний нужно~--- это и есть первый из шести шагов
ниже.
\end{buildbox}

\section{Шесть шагов проектирования}

\begin{enumerate}
  \item \textbf{Словесное описание.} Что должно делать устройство, какие у него
        входы и выходы, когда выход должен меняться.
  \item \textbf{Граф состояний.} Нарисовать кружки (состояния) и стрелки
        (переходы) с подписями «условие / выход».
  \item \textbf{Таблица переходов.} Для каждого состояния и каждого значения
        входов записать следующее состояние и выход. Ни одна ячейка не должна
        остаться пустой.
  \item \textbf{Кодирование состояний.} Назначить состояниям двоичные коды
        или использовать унитарный код.
  \item \textbf{Логика.} Получить функции следующего состояния и выходов
        (вручную по картам Карно или с помощью программы).
  \item \textbf{Проверка.} Перебрать сценарии на симуляции, включая
        граничные: сброс в середине работы, лишние входы, одновременные события.
\end{enumerate}

\begin{diybox}[Шесть шагов на маленьком примере]
\textbf{Задание.} Пройти все шесть шагов для автомата «лампа с кнопкой»: одно
нажатие включает, следующее выключает.\\
\textbf{В Logra.} Вкладка \emph{Обучение}, книга «из блоков», глава
«Автоматы»; вкладки \emph{Конструктор} и \emph{Диаграммы}.\\
\stars{1} \emph{Повтори.} Урок «Автомат на двух состояниях».\\
\stars{2} \emph{Измени.} Урок «Каркас автомата»: добавьте третье состояние
«мигает», в которое лампа переходит после двойного нажатия.\\
\stars{3} \emph{Проект.} Светофор для пешеходов (урок «Светофор для
пешеходов»): нарисуйте граф и таблицу переходов \emph{до} сборки; при
сборке ни одна ячейка таблицы не должна остаться пустой.\\
\textbf{Ожидаемый результат.} Диаграмма повторяет путь по графу: состояния
сменяются в том порядке, который вы нарисовали.\\
\textbf{Как проверить.} Вкладка \emph{Диаграммы}: сверьте каждый такт с
таблицей переходов; вкладка \emph{RTL}: найдите регистр состояния и логику
перехода.
\end{diybox}

\section{Пример: автомат по продаже}

Автомат продаёт товар за 15~рублей. Он принимает монеты по 5 и по 10 рублей,
по одной за такт. Когда внесено 15 рублей или больше, он \emph{в тот же
такт} выдаёт сигнал \code{vend}, после чего начинает считать заново (сдачу
для простоты не выдаём).

Состояний три: внесено 0, 5 или 10 рублей. Выход зависит не только от
состояния, но и от монеты, которая пришла сейчас: пятёрка в состоянии
«10 рублей» сразу даёт продажу. Значит, перед нами автомат \term{Мили}. Его
граф на рис.~\ref{fig:vend_graph}: на стрелках подписаны входы и выход через
косую черту.

\begin{figure}[H]
\centering
\begin{tikzpicture}[scale=0.85, transform shape, font=\sffamily\small, st/.style={circle, draw=FpgaBlue, fill=FpgaLight, minimum size=16mm, thick, align=center}, lbl2/.style={font=\sffamily\scriptsize, fill=white, inner sep=1.5pt}]
  \node[st] (s0) at (0,0) {$S_0$\\{\scriptsize 0 руб.}};
  \node[st] (s5) at (5,2.6) {$S_5$\\{\scriptsize 5 руб.}};
  \node[st] (s10) at (10,0) {$S_{10}$\\{\scriptsize 10 руб.}};
  \draw[-{Stealth}] (s0) -- node[lbl2, above left] {c5/0} (s5);
  \draw[-{Stealth}] (s5) -- node[lbl2, above right] {c5/0} (s10);
  \draw[-{Stealth}] (s5) to[bend right=30] node[lbl2, below left] {c10/1} (s0);
  \draw[-{Stealth}] (s0) to[bend left=18] node[lbl2, above] {c10/0} (s10);
  \draw[-{Stealth}] (s10) to[bend left=18] node[lbl2, below] {c5 или c10/1} (s0);
  \draw[-{Stealth}] (s0) to[loop left] node[lbl2, left] {нет монеты/0} (s0);
  \draw[-{Stealth}] (s5) to[loop above] node[lbl2, above] {нет монеты/0} (s5);
  \draw[-{Stealth}] (s10) to[loop right] node[lbl2, right] {нет монеты/0} (s10);
\end{tikzpicture}
\caption{Граф автомата Мили. Надпись «вход/выход»: после косой черты значение \code{vend}. Все переходы, которые заканчиваются продажей, возвращают автомат в $S_0$}
\label{fig:vend_graph}
\end{figure}

Те же сведения собраны в таблице переходов. В ней не должно быть пустых
ячеек:

\begin{table}[H]
\centering
\caption{Таблица переходов автомата по продаже ($c5$, $c10$ --- монеты, \code{vend} --- выход)}
\small
\begin{tabular}{lcccc}
\toprule
\textbf{Состояние} & \textbf{Нет монеты} & \textbf{c5} & \textbf{c10} & \textbf{Выход} \\
\midrule
$S_0$ (0) & $S_0$ & $S_5$ & $S_{10}$ & 0 \\
$S_5$ (5) & $S_5$ & $S_{10}$ & $S_0$ & 1 только при c10 \\
$S_{10}$ (10) & $S_{10}$ & $S_0$ & $S_0$ & 1 при c5 и при c10 \\
\bottomrule
\end{tabular}
\end{table}

\subsection{Как автомат работает такт за тактом}

Проверим автомат вручную: подадим монеты $5, 5, 5, 10, 5$ и выпишем в таблицу,
что происходит на каждом такте. Состояние и монета определяют и выход, и
следующее состояние.

\begin{table}[H]
\centering
\caption{Трасса работы автомата по продаже}
\small
\begin{tabular}{cllcl}
\toprule
\textbf{Такт} & \textbf{Состояние} & \textbf{Монета} & \textbf{\code{vend}} & \textbf{Следующее состояние} \\
\midrule
0 & $S_0$ & 5 & 0 & $S_5$ \\
1 & $S_5$ & 5 & 0 & $S_{10}$ \\
2 & $S_{10}$ & 5 & \textbf{1} & $S_0$ \\
3 & $S_0$ & 10 & 0 & $S_{10}$ \\
4 & $S_{10}$ & 5 & \textbf{1} & $S_0$ \\
5 & $S_0$ & нет & 0 & $S_0$ \\
\bottomrule
\end{tabular}
\end{table}

Обратите внимание на такты 2 и 4: сигнал \code{vend} равен единице \emph{в том
же такте}, в котором брошена последняя монета, а не на следующем. Выход
вычисляется по состоянию и по входу вместе, и это главная примета автомата
Мили. Автомат Мура выдал бы товар только в следующем такте, когда регистр
состояния уже сменил значение.

\subsection{Правила надёжного описания}

Каждый автомат делят на три части: регистр состояния, комбинационную логику
перехода и комбинационную логику выхода. Три приёма защищают от типичных
ошибок.

\begin{itemize}
  \item Для каждого состояния заранее определено «следующее состояние по
        умолчанию = текущее»: если ни одно условие перехода не сработало,
        автомат остаётся на месте, и в схеме не появляется лишней защёлки.
  \item Для неиспользуемых кодов состояния предусматривают переход «во всех
        остальных случаях»: если регистр каким-то образом попадёт в такой код
        (например, из-за сбоя), автомат вернётся в начальное состояние.
  \item Сброс делают синхронным: он действует по такту и поэтому не создаёт
        проблем с временем.
\end{itemize}

Как эти три приёма выглядят в коде и как этот автомат проверяют на
последовательности монет, показано в разделе~\ref{sec:v-fsm3}.

Сравним два вида автоматов.

\begin{table}[H]
\centering
\caption{Автоматы Мура и Мили}
\small
\begin{tabularx}{\textwidth}{l X X}
\toprule
 & \textbf{Мур} & \textbf{Мили} \\
\midrule
Выход зависит от & состояния & состояния и входов \\
Число состояний & обычно больше & обычно меньше \\
Реакция на вход & через такт & сразу, в том же такте \\
Помехи на входе & не доходят до выхода & доходят до выхода \\
\bottomrule
\end{tabularx}
\end{table}

\begin{diybox}[Автомат по продаже]
\textbf{Задание.} Проверить автомат по продаже на последовательностях монет
и научить его обрабатывать одновременные монеты.\\
\textbf{В Logra.} Вкладка \emph{Примеры}, проект \code{vending-fsm}
(автомат считает события и выдаёт результат после заданного числа);
вкладки \emph{Диаграммы} и \emph{RTL}.\\
\stars{1} \emph{Повтори.} Сравните одно длинное удержание кнопки и три
отдельных нажатия с отпусканиями: на диаграмме найдите, когда появляется
импульс выдачи.\\
\stars{2} \emph{Измени.} На вкладке \emph{Код} увеличьте число событий до
выдачи с трёх до четырёх и проверьте импульс выдачи на диаграмме. Затем
составьте по таблице главы трассу для монет $10, 5, 10$.\\
\stars{3} \emph{Проект.} Задача~2 главы: добавьте в таблицу переходов столбец
«обе монеты сразу» и проверьте таблицу на пустые ячейки.\\
\textbf{Ожидаемый результат.} Состояние накапливает события, а после
заданного числа выдаёт однотактовый импульс и возвращается в начало.\\
\textbf{Как проверить.} Вкладка \emph{Диаграммы}: между импульсами выдачи
ровно столько событий, сколько вы задали; вкладка \emph{RTL}: найдите регистр
накопленной суммы.
\end{diybox}

\section{Автомат в три блока на Verilog}
\label{sec:v-fsm3}

Граф и таблица переходов переводятся в Verilog почти механически. Автомат
записывают тремя отдельными частями: каждая отвечает на свой вопрос, и их
нельзя смешивать, иначе в схеме появляются защёлки и трудноуловимые ошибки.

\begin{table}[H]
\centering
\caption{Три блока автомата и их запись}
\small
\begin{tabularx}{\textwidth}{l X X}
\toprule
\textbf{Блок} & \textbf{Что делает} & \textbf{Как пишется} \\
\midrule
1. Регистр состояния & запоминает, где мы сейчас; сброс ставит начальное состояние & \code{always @(posedge clk)}, присваивание \code{<=}, сброс первым \\
2. Логика переходов & по состоянию и входам считает следующее состояние \code{next} & \code{always @(*)}, \code{next = state;} в начале, затем \code{case (state)} с веткой \code{default} \\
3. Логика выходов & Мур: выход зависит только от \code{state}. Мили: от \code{state} и входов & \code{always @(*)} с \code{case} или \code{assign}; присваивание \code{=} \\
\bottomrule
\end{tabularx}
\end{table}

\subsection{Светофор: автомат Мура}

Три состояния, в каждом горит одна лампа; время в состоянии отсчитывает
счётчик тактов \code{t}. Длительности вынесены в параметры, поэтому
изменить зелёный с 4 тактов на 6 можно одним числом.

\vfile{\codepath/b9a_fsm/traffic3.v}{code/b9a\_fsm/traffic3.v}

\textbf{Как читается.} Блок~1~--- единственное место, где меняются
\code{state} и \code{t}: по фронту такта состояние берёт значение
\code{next}. В блоке~2 первая строка \code{next = state} задаёт «остаться»
по умолчанию, а сигнал \code{last} («это последний такт состояния») решает,
уходить ли. Ветка \code{default} возвращает автомат в начало, если регистр
попал в неиспользуемый код \code{3}. Блок~3 зависит только от \code{state}:
это и есть автомат Мура, поэтому лампа меняется ровно вместе с состоянием.

\shot{br-fsm-wave}{Светофор на диаграмме: на каждом состоянии горит одна лампа ($0x1$ --- зелёная, $0x2$ --- жёлтая, $0x4$ --- красная); снимок модуля \code{traffic\_fsm} из \code{code/00\_theory}, у \code{traffic3} картина та же}{fig:b-fsm-wave}

Тестбенч \code{b9a\_fsm/tb\_traffic3.v} проверяет три круга по формуле:
зелёный 4 такта, жёлтый 2, красный 4. Он печатает \code{traffic3: OK}.

\begin{diybox}[Светофор по параметрам]
\textbf{Задание.} Запустить светофор и поменять длительности.\\
\textbf{В Logra.} Вкладка «Код»: вставьте \code{traffic3.v}, «Синтезировать»
(Ctrl+Enter), вкладки RTL и «Диаграммы» (число циклов 30, импульс \code{rst} в
начале).\\
\textbf{Ожидаемый результат.} \code{lamp} принимает значения \code{001},
\code{010}, \code{100} в течение 4, 2 и 4 тактов и повторяется.\\
\textbf{Как проверить.} На RTL найдите регистр \code{state}, регистр \code{t}
и мультиплексор выхода. В терминале:
\code{iverilog -g2012 -o tb tb\_traffic3.v traffic3.v}, затем \code{vvp tb}.
\stars{2} Добавьте четвёртое состояние: перед зелёным на 1 такт горят красный и
жёлтый вместе (\code{3'b110}).\\
\textbf{Образец ответа.} Нужны новая константа, ветка в блоках~2 и~3 и
изменение перехода из красного.
\end{diybox}

\begin{vcode}
module traffic4 (
    input            clk, rst,
    output reg [2:0] lamp
);
    localparam [1:0] RED = 2'd0, RED_YEL = 2'd1, GREEN = 2'd2, YELLOW = 2'd3;
    reg [1:0] state, next;
    reg [2:0] t;
    wire last = (state == RED && t == 3'd3) || (state == RED_YEL && t == 3'd0)
             || (state == GREEN && t == 3'd3) || (state == YELLOW && t == 3'd1);
    always @(posedge clk) begin
        if (rst) begin state <= RED; t <= 3'd0; end
        else begin state <= next; t <= last ? 3'd0 : t + 3'd1; end
    end
    always @(*) begin
        next = state;
        if (last) case (state)
            RED:     next = RED_YEL;
            RED_YEL: next = GREEN;
            GREEN:   next = YELLOW;
            default: next = RED;
        endcase
    end
    always @(*) begin
        case (state)
            RED:     lamp = 3'b100;
            RED_YEL: lamp = 3'b110;
            GREEN:   lamp = 3'b001;
            default: lamp = 3'b010;
        endcase
    end
endmodule
\end{vcode}


\subsection{Автомат по продаже: Мили}

Автомат из начала главы записывается теми же тремя блоками. Разница только
в блоке~3: выход \code{vend} вычисляется по состоянию \emph{и по входам}
\code{c5}, \code{c10}.

\vfile{\codepath/00_theory/vending.v}{code/00\_theory/vending.v}

\textbf{Как читается.} Состояния~--- сколько рублей внесено. В блоке~2 для
каждого состояния описаны обе монеты, а при отсутствии монеты остаётся
\code{next = state}. Блок~3~--- одна формула: продажа, если в $S_5$ пришла
десятка или в $S_{10}$ пришла любая монета. Три приёма из раздела «Правила
надёжного описания» все на месте: значение по умолчанию, \code{default} и
синхронный сброс.

\shot{n-vend-wave}{Автомат Мили: \code{vend} вырастает в том же такте, в котором приходит последняя монета}{fig:n-vend-wave}

Монеты на «Диаграммах» задают в разделе «События». Сигнал \code{vend}
поднимается внутри такта, как только на входе появилась «решающая» монета, не
дожидаясь следующего фронта~--- в этом отличие автомата Мили. Если нужен
«чистый» выход, его пропускают через триггер, и он запаздывает на такт.
Тестбенч \code{00\_theory/tb\_vending.v} бросает те же монеты, что трасса в
таблице главы (5+5+5, 10+5, 5+10, 10+10), и считает продажи: должно быть четыре.

\begin{diybox}[Автомат по продаже в коде]
\textbf{Задание.} Запустить автомат и научить его новому правилу.\\
\textbf{В Logra.} Вкладка «Код»: \code{vending.v}, «Синтезировать», «Диаграммы»:
раздел «События», импульсы \code{c5} и \code{c10} по тактам.\\
\textbf{Ожидаемый результат.} Для монет 5, 5, 5 \code{vend = 1} в такте третьей
монеты; для монет 10, 5~--- в такте пятёрки.\\
\textbf{Как проверить.} Запустите \code{tb\_vending.v}: \code{vending: OK}.
\stars{2} Сделайте, чтобы цена была 20 рублей (добавьте состояние $S_{15}$).
\end{diybox}

\textbf{Образец ответа.}
\begin{vcode}
module vending20 (
    input  clk, rst, c5, c10,
    output vend
);
    localparam [1:0] S0 = 2'd0, S5 = 2'd1, S10 = 2'd2, S15 = 2'd3;
    reg [1:0] state, next;
    always @(posedge clk) begin
        if (rst) state <= S0;
        else     state <= next;
    end
    always @(*) begin
        next = state;
        case (state)
            S0:  if (c10) next = S10; else if (c5) next = S5;
            S5:  if (c10) next = S15; else if (c5) next = S10;
            S10: if (c10) next = S0;  else if (c5) next = S15;
            S15: if (c10 | c5) next = S0;
        endcase
    end
    assign vend = (state == S10 && c10) || (state == S15 && (c5 | c10));
endmodule
\end{vcode}


\subsection{Детектор последовательности 1-0-1}

Детектор из главы~\ref{ch:sync} и из кода панели~--- тоже автомат Мили.
Состояния: \code{S0}~--- «ничего полезного», \code{S1}~--- «видели 1»,
\code{S10}~--- «видели 10». Выход \code{hit} равен единице, когда в состоянии
\code{S10} приходит ещё одна единица. Единица в любом состоянии может начать
новую последовательность, поэтому тройки перекрываются: поток \code{10101}
даёт два срабатывания.

\vfile{\codepath/b9a_fsm/seq101.v}{code/b9a\_fsm/seq101.v}

\textbf{Как читается.} Блок~2 содержит всю «память» последовательности: из
\code{S10} при входе 1 автомат переходит не в \code{S0}, а в \code{S1}, потому
что только что пришедшая единица уже начала новую тройку. Блок~3~--- одна
строка, выход Мили.

Тестбенч \code{tb\_seq101.v} подаёт поток \code{1 0 1 0 1 1 0 1}, где три
срабатывания, и поток \code{1 1 0 0 0 1 1 1}, где ни одного. Подача бита
\code{send(b)} смотрит на \code{hit} \emph{до} фронта такта: так проверяют
выход Мили.

\begin{diybox}[Свой детектор]
\textbf{Задание.} Запустить детектор и переделать его на другой код.\\
\textbf{В Logra.} Вкладка «Код»: \code{seq101.v}; «Диаграммы»: раздел «События»,
сигнал \code{x} по тактам \code{1, 0, 1, 0, 1}. Либо урок «Кодовый замок»
(вкладка «Обучение», книга по Verilog).\\
\textbf{Ожидаемый результат.} Два импульса \code{hit} на потоке \code{10101};
ни одного на \code{110} и \code{011}.\\
\textbf{Как проверить.} \code{tb\_seq101.v} печатает \code{seq101: OK}.
\stars{2} Переделайте детектор на код \code{1-1-0}.\\
\textbf{Образец ответа.} Меняются переходы и выход: «видели 1», «видели 11» и
срабатывание на нуле после двух единиц.
\end{diybox}

\begin{vcode}
module seq110 (
    input  clk, rst, x,
    output hit
);
    localparam [1:0] S0 = 2'd0, S1 = 2'd1, S11 = 2'd2;
    reg [1:0] state, next;
    always @(posedge clk) begin
        if (rst) state <= S0;
        else     state <= next;
    end
    always @(*) begin
        next = state;
        case (state)
            S0:      next = x ? S1 : S0;
            S1:      next = x ? S11 : S0;
            S11:     next = x ? S11 : S0;   // 1-1-1 ещё содержит 1-1
            default: next = S0;
        endcase
    end
    assign hit = (state == S11) && !x;
endmodule
\end{vcode}


\section{Кодирование и сокращение состояний}

Три состояния требуют двух триггеров при двоичном кодировании
(\code{00}, \code{01}, \code{10}) и трёх --- при унитарном. Выбор кодирования
может изменить размер и скорость схемы. Синтезатор (глава~\ref{ch:flow})
обычно может перекодировать автомат сам, если описать его по шаблону с
именованными состояниями; но выбор кодирования полезно уметь сделать и
вручную.

Иногда два состояния ведут себя одинаково: при любых входах они дают те же
выходы и переходят в одинаковые состояния. Такие состояния можно слить в
одно, и автомат станет проще. Проверка проста: если у двух состояний
совпадают все строки таблицы переходов (с точностью до имён этих же
состояний), их можно объединить.

Покажем три кодирования на одном автомате. Четыре состояния \code{A}, \code{B},
\code{C}, \code{D} идут по кругу, переход делается по сигналу \code{go}, а
выход \code{y} равен единице в состояниях \code{C} и \code{D}. В двоичном
коде их записывают числами 0--3, в коде Грея соседние состояния отличаются
одним битом (\code{00, 01, 11, 10}), в унитарном каждому состоянию отвечает
свой триггер и горит ровно один.

\vfile{\codepath/b9a_fsm/state_enc.v}{code/b9a\_fsm/state\_enc.v}

\textbf{Как читается.} В первых двух модулях одинаковый \code{case}, различаются
только константы \code{A\ldots D}: выбор кода ничего не меняет в тексте
автомата. Зато в Грее выход оказался просто старшим битом \code{st[1]}, и
выходная логика исчезла. В унитарном модуле нет \code{case}: переход~--- это
сдвиг единицы на соседний триггер, а выход~--- ИЛИ двух триггеров. Цена:
4 триггера вместо 2, выигрыш: почти нет логики между триггерами. Тестбенч
\code{tb\_state\_enc.v} подаёт одни и те же \code{go} на все три модуля и
убеждается, что выходы совпадают на каждом такте.

\begin{diybox}[Двоичное и унитарное кодирование]
\textbf{Задание.} Построить один и тот же автомат двумя способами и сравнить
цену каждого.\\
\textbf{В Logra.} Вкладки \emph{Конструктор}, \emph{Диаграммы}, \emph{RTL},
\emph{Gowin}.\\
\stars{1} \emph{Повтори.} Для автомата «красный~--- зелёный~--- жёлтый»
выпишите двоичные коды состояний и унитарные.\\
\stars{2} \emph{Измени.} Соберите автомат в Конструкторе на двух триггерах
(двоичное кодирование), затем на трёх (унитарное); диаграммы выходов должны
совпасть.\\
\stars{3} \emph{Проект.} Сравните число триггеров и логических элементов на
вкладках \emph{RTL} и \emph{Gowin}; запишите вывод в одну строку.\\
\textbf{Ожидаемый результат.} Унитарный вариант использует больше триггеров
и меньше логики на каждый переход.\\
\textbf{Как проверить.} Вкладка \emph{Диаграммы}: выходы обоих вариантов
идентичны на всех тактах; вкладка \emph{RTL}: подсчёт триггеров.\\
\textbf{То же на Verilog.} Вкладка \emph{Код}: вставьте \code{state\_enc.v},
«Синтезировать», по очереди выберите верхним \code{ring\_bin},
\code{ring\_gray}, \code{ring\_onehot} и сравните число регистров на RTL (2, 2
и 4). \stars{2} Допишите унитарное кольцо на три состояния светофора.
\end{diybox}

\textbf{Образец ответа.}
\begin{vcode}
module ring3_onehot (
    input      clk, rst, go,
    output     green, yellow, red
);
    reg [2:0] st;                       // st[0] = зелёный, st[1] = жёлтый, st[2] = красный
    always @(posedge clk) begin
        if (rst)     st <= 3'b001;
        else if (go) st <= {st[1:0], st[2]};
    end
    assign green  = st[0];
    assign yellow = st[1];
    assign red    = st[2];
endmodule
\end{vcode}


\section{Сквозной проект: автомат «Охранная панель»}
\label{sec:panel-fsm}

До сих пор в проекте были отдельные узлы: датчик, логика «тревога», запоминание
тревоги, счётчик времени. Теперь соберём из них устройство с режимами. Панель
находится в одном из трёх состояний: \term{ОХРАНА\_ВЫКЛ} (датчики
игнорируются, жильцы дома), \term{ОХРАНА} (датчик двери следит) и
\term{ТРЕВОГА} (включена сирена). Переключают режимы вводом кода: на кнопку
по тактам подаются биты, и код считается введённым, если подряд пришло
\code{1-0-1}.

Здесь два автомата соединены одним сигналом. Первый~--- детектор кода из
главы~\ref{ch:sync}: он следит за вводом и выдаёт $c=1$ в том такте, когда код
набран верно. Второй~--- автомат режимов, который читает сигнал $c$ и датчик
двери $d$ (1~--- дверь открыта).

\begin{figure}[H]
\centering
\begin{tikzpicture}[
  st/.style={circle, draw=FpgaDark, very thick, minimum size=19mm, font=\sffamily\small, align=center, fill=FpgaLight},
  tr/.style={-{Stealth[length=3mm]}, thick, FpgaDark},
  lb/.style={font=\sffamily\scriptsize, fill=white, inner sep=1.5pt}]
  \node[st] (off) at (0,0) {ОХРАНА\\ВЫКЛ\\\scriptsize сирена 0};
  \node[st] (arm) at (5.2,0) {ОХРАНА\\\scriptsize сирена 0};
  \node[st, fill=red!12] (al) at (10.4,0) {ТРЕВОГА\\\scriptsize сирена 1};
  \draw[tr] (off) to[bend left=22] node[lb, above] {$c=1$} (arm);
  \draw[tr] (arm) to[bend left=22] node[lb, below] {$c=1$} (off);
  \draw[tr] (arm) -- node[lb, above] {$d=1$, $c=0$} (al);
  \draw[tr] (al) to[bend left=50] node[lb, below] {$c=1$} (off);
  \draw[tr] (off) edge[loop above, looseness=7] node[lb, above] {$c=0$} (off);
  \draw[tr] (arm) edge[loop above, looseness=7] node[lb, above] {$d=0$, $c=0$} (arm);
  \draw[tr] (al) edge[loop above, looseness=7] node[lb, above] {$c=0$} (al);
  \draw[tr, FpgaRed] (-1.9,-1.9) node[left, font=\sffamily\scriptsize, text=FpgaRed] {сброс} -- (off.south west);
\end{tikzpicture}
\caption{Автомат режимов панели (Мура). Сирена включена только в состоянии
ТРЕВОГА. Условие выхода из любого режима в ОХРАНА\_ВЫКЛ~--- верный код}
\label{fig:panel_fsm}
\end{figure}

Проверим по шести шагам. Состояния заданы (шаги 1--2). Таблица переходов
(шаг 3) не имеет пустых ячеек: для каждого состояния описаны все четыре
комбинации $(d,c)$. В состоянии ОХРАНА при $d=1$ и $c=1$ побеждает код:
хозяин вошёл и ввёл код, сирена не нужна.

\begin{table}[H]
\centering
\caption{Таблица переходов автомата панели ($d$~--- дверь открыта, $c$~--- верный код)}
\small
\begin{tabular}{l|cccc|c}
\toprule
\textbf{Состояние} & $d{=}0,c{=}0$ & $d{=}0,c{=}1$ & $d{=}1,c{=}0$ & $d{=}1,c{=}1$ & \textbf{Сирена} \\
\midrule
ОХРАНА\_ВЫКЛ & ВЫКЛ & ОХРАНА & ВЫКЛ & ОХРАНА & 0 \\
ОХРАНА       & ОХРАНА & ВЫКЛ & ТРЕВОГА & ВЫКЛ & 0 \\
ТРЕВОГА      & ТРЕВОГА & ВЫКЛ & ТРЕВОГА & ВЫКЛ & 1 \\
\bottomrule
\end{tabular}
\end{table}

Теперь запишем оба автомата в Verilog. Детектор кода~--- тот же \code{seq101}
(в файле он переименован в \code{code\_detector}, чтобы файл был
самодостаточным). Автомат режимов~--- три блока: регистр \code{mode}, переходы
по таблице выше и выход Мура \code{siren}. Верхний модуль \code{panel\_fsm}
только соединяет их сигналом \code{c}.

\vfile{\codepath/b9a_fsm/panel_fsm.v}{code/b9a\_fsm/panel\_fsm.v}

\textbf{Как читается.} В блоке переходов ветка \code{ARMED} сначала проверяет
\code{c}, и лишь потом \code{d}: так в клетке $d{=}1,c{=}1$ побеждает код,
как требует таблица. Состояние \code{ALARM} выходит только по коду: тревога
сама не гаснет, даже если дверь закрыли. Тестбенч \code{tb\_panel\_fsm.v}
проходит четыре пункта критерия приёмки из «Сделай сам» ниже и печатает
\code{panel\_fsm: OK}.

\begin{diybox}[Этап 9 проекта: автомат режимов]
\textbf{Задание.} Собрать автомат «Охранная панель» с вводом кода \code{1-0-1}
и сиреной.\\
\textbf{В Logra.} Вкладка \emph{Обучение}, книга «из блоков», глава
«Автоматы»: урок «Кодовый замок» (код \code{1, 1, 0}); вкладка \emph{Код},
книга по Verilog: урок «Кодовый замок: 1, 0, 1»; вкладка \emph{Примеры}:
проект \code{door-controller} как образец автомата с режимами; вкладки
\emph{Конструктор}, \emph{Диаграммы}, \emph{RTL}.\\
\stars{1} \emph{Повтори.} Урок «Кодовый замок» из блоков: соберите замок и
убедитесь, что он открывается только нужной тройкой.\\
\stars{2} \emph{Измени.} Перестройте детектор под код \code{1-0-1}, затем
добавьте в автомат режимов вход $d$ и выход «сирена» по таблице главы.\\
\stars{3} \emph{Проект.} \textbf{Этап проекта «Охранная панель»:} автомат
ОХРАНА\_ВЫКЛ $\to$ ОХРАНА $\to$ ТРЕВОГА с вводом кода \code{1-0-1} и сбросом
тревоги тем же кодом. Критерий приёмки: на диаграмме (а) после ввода \code{101}
панель переходит в ОХРАНА, (б) открытие двери ($d=1$) в режиме ОХРАНА включает
сирену на следующем такте, (в) в режиме ОХРАНА\_ВЫКЛ открытие двери сирену
\emph{не} включает, (г) ввод \code{101} в состоянии ТРЕВОГА выключает сирену.\\
\textbf{Ожидаемый результат.} Все четыре пункта критерия выполняются на
диаграмме, таблица переходов не имеет пустых ячеек.\\
\textbf{Как проверить.} Вкладка \emph{Диаграммы}: по очереди подайте код,
откройте дверь, снова введите код; вкладка \emph{RTL}: найдите регистр
состояния и две логики (переходов и выхода).\\
\textbf{То же на Verilog.} Вкладка \emph{Код}: вставьте \code{panel\_fsm.v},
«Синтезировать», на «Диаграммах» подайте \code{code\_bit} по тактам
\code{1, 0, 1}, затем \code{d = 1}. Образец ответа к \stars{3}~--- сам файл
\code{panel\_fsm.v}; тестбенч \code{tb\_panel\_fsm.v} проверяет пункты (а)--(г).
\end{diybox}

\section{Очередь FIFO}

Часто источник данных и их потребитель работают с разной скоростью: датчик
отдаёт отсчёты пачками, а линия связи отправляет их медленно. Нужен буфер, который
запоминает данные и выдаёт их в том же порядке. Такой буфер называют
\term{FIFO} (first in~--- first out, «первым пришёл --- первым ушёл»). Работает
он как очередь в магазине.

\subsection{Кольцевой буфер}

Внутри FIFO стоит обычная память (массив слов) и два указателя:
\emph{указатель записи} (куда класть следующее слово) и \emph{указатель
чтения} (откуда брать). Каждая операция сдвигает свой указатель на единицу, а
в конце памяти он возвращается к нулю, как стрелка на циферблате. Слов в
очереди столько, сколько накопилось между указателями: когда оно равно нулю,
очередь пуста (\code{empty}), когда равно размеру памяти --- заполнена
(\code{full}).

\begin{figure}[H]
\centering
\begin{tikzpicture}[font=\sffamily\small]
  \foreach \i in {0,...,7} {
    \pgfmathsetmacro{\a}{90-\i*45}
    \node[draw, minimum size=9mm, fill=white] (c\i) at (\a:2.2) {\i};
  }
  % занятые ячейки: 2,3,4
  \foreach \i in {2,3,4} \node[draw, minimum size=9mm, fill=FpgaLight] at (90-\i*45:2.2) {\i};
  \node[font=\sffamily\bfseries, text=FpgaGreen!50!black] at (90-2*45:3.5) {rp};
  \draw[-{Stealth}, FpgaGreen!50!black, thick] (90-2*45:3.3) -- (90-2*45:2.8);
  \node[font=\sffamily\bfseries, text=FpgaOrange] at (90-5*45:3.5) {wp};
  \draw[-{Stealth}, FpgaOrange, thick] (90-5*45:3.3) -- (90-5*45:2.8);
  \node[align=left, anchor=west] at (4.4,1.5) {Хранится 3 слова (ячейки 2, 3, 4).\\Читать начнём с ячейки 2.\\Следующая запись пойдёт в ячейку 5.};
  \node[align=left, anchor=west] at (4.4,-1.2) {После ячейки 7 указатели\\возвращаются к 0: поэтому для\\8 слов хватает 3-разрядных\\счётчиков.};
\end{tikzpicture}
\caption{Кольцевой буфер: заполнены ячейки между указателем чтения \code{rp} и указателем записи \code{wp}}
\end{figure}

\subsection{Признаки full и empty}

Для очереди на 8 слов хватает трёхразрядных указателей, потому что
счётчик из трёх бит сам «заворачивается» после значения 7. Отдельный
четырёхразрядный счётчик слов \code{count} (значения от 0 до 8) упрощает
признаки \code{empty} и \code{full}. Запись разрешена, только пока очередь не
полна, а чтение --- пока она не пуста.

\begin{table}[H]
\centering
\caption{Трасса очереди: записали три слова, затем прочитали три}
\small
\begin{tabular}{lcccc}
\toprule
\textbf{Действие} & \textbf{wp} & \textbf{rp} & \textbf{count} & \textbf{empty} \\
\midrule
после сброса & 0 & 0 & 0 & 1 \\
записали 11 & 1 & 0 & 1 & 0 \\
записали 22 & 2 & 0 & 2 & 0 \\
записали 33 & 3 & 0 & 3 & 0 \\
прочитали (получили 11) & 3 & 1 & 2 & 0 \\
прочитали (получили 22) & 3 & 2 & 1 & 0 \\
прочитали (получили 33) & 3 & 3 & 0 & 1 \\
\bottomrule
\end{tabular}
\end{table}

Слова выходят в том порядке, в котором вошли. Заметьте, что равенство
указателей $\text{wp}=\text{rp}$ не отличает пустую очередь от полной: в обоих
случаях указатели совпадают. Поэтому признаки берут не из сравнения
указателей, а из счётчика слов: $\text{empty}=(\text{count}=0)$,
$\text{full}=(\text{count}=8)$.

\subsection{Переполнение и опустошение на диаграммах}

Что произойдёт, если источник пишет, а очередь уже полна, или приёмник читает
из пустой очереди? Рассмотрим очередь на 4 слова. Пять тактов подряд
пытаемся писать, затем пять тактов подряд читать.

\begin{figure}[H]
\centering
\begin{tikztimingtable}[timing/slope=0.05, xscale=1.1, yscale=1.4, timing/font=\sffamily\small, timing/rowdist=3.2ex]
  CLK    & 10{1L 1H} \\
  wr     & 10H 10L \\
  rd     & 10L 10H \\
  count  & 2D{0} 2D{1} 2D{2} 2D{3} 2D{4} 2D{4} 2D{3} 2D{2} 2D{1} 2D{0} \\
  full   & 8L 4H 8L \\
  empty  & 2H 16L 2H \\
\end{tikztimingtable}
\caption{Очередь на 4 слова. На пятом такте очередь уже полна (\code{full}=1): пятая запись
игнорируется, слово теряется. На последнем такте очередь пуста (\code{empty}=1): десятое
чтение ничего не выдаёт, счётчик остаётся нулевым}
\label{fig:fifo_flags}
\end{figure}

Очередь «защищает» себя: при $\text{full}=1$ запись не сдвигает указатель, а
при $\text{empty}=1$ чтение не сдвигает свой. Но слово, пришедшее на
полную очередь, \emph{потеряно}: исправить это может только тот, кто пишет,
если будет смотреть на \code{full} и ждать. Так что в системе источник
обязан либо проверять признак, либо иметь запас по размеру очереди.

\subsection{Пример: очередь на четыре слова}

\begin{buildbox}[Посмотрите очередь в работе]
Вкладка \emph{Примеры}, проект \code{fifo-four} («FIFO на четыре слова»).
После сброса внутренний источник сам записывает слова 1, 2, 3, 4 и
останавливается, когда очередь полна. Каждое нажатие кнопки читает самое
старое слово. Указатели \code{write\_ptr} и \code{read\_ptr} двигаются
независимо, а счётчик \code{used} ограничивает доступ. Проследите на
\emph{Диаграммах}, как указатели «оборачиваются».
\end{buildbox}

\subsection{Очередь на Verilog: счётчик слов}

Тот же буфер на восемь слов по восемь бит. Массив \code{mem}~--- память,
\code{wp} и \code{rp}~--- указатели записи и чтения, а \code{count}~--- счётчик
слов, по которому берутся признаки.

\vfile{\codepath/00_theory/fifo8.v}{code/00\_theory/fifo8.v}

\textbf{Как читается.} Условия \code{do\_wr} и \code{do\_rd} включают защиту:
запись в полную очередь и чтение из пустой не меняют ничего. Трёхразрядные
\code{wp} и \code{rp} сами «заворачиваются» после семи, поэтому отдельной
проверки на конец памяти не нужно. Строка
\code{count <= count + do\_wr - do\_rd} учитывает одновременные запись и
чтение: счётчик не меняется. Массив \code{mem} синтезатор размещает в
памяти кристалла или в ячейках.

\shot{n-fifo-wave}{Работа очереди: записаны слова \code{0B}, \code{16}, \code{21}, затем прочитаны в том же порядке}{fig:n-fifo-wave}

На диаграмме \code{count} растёт с каждой записью и убывает с каждым чтением,
\code{empty} сбрасывается после первой записи, а слова выходят на \code{dout} в
порядке записи. Тестбенч \code{00\_theory/tb\_fifo8.v} пишет десять слов (два
последних не помещаются), проверяет, что \code{full = 1}, читает восемь слов по
порядку и убеждается, что очередь пуста.

\subsection{Очередь на Verilog: лишний бит указателя}

Счётчик слов можно не заводить. Если сделать указатели \emph{на один бит
шире} адреса, то старший бит считает «круги»: очередь пуста, когда указатели
равны целиком, и полна, когда равны младшие биты (адрес), а старшие различаются
(писатель обогнал читателя ровно на круг). Так пишут очереди с параметрами:
глубина $2^{\text{AW}}$, ширина слова \code{DW}.

\vfile{\codepath/b9a_fsm/fifo_ptr.v}{code/b9a\_fsm/fifo\_ptr.v}

\textbf{Как читается.} Адрес в память берётся из \code{wp[AW-1:0]} и
\code{rp[AW-1:0]}, а признаки~--- из полных указателей. Слово \code{dout}
показывается сразу (без ожидания такта): оно всегда равно самому старому
слову, пока очередь не пуста. Тестбенч \code{tb\_fifo\_ptr.v} сравнивает очередь
с простой моделью (массив и два целых индекса) на 12 записях подряд, 12
чтениях и 400 случайных тактах: проверяются и \code{full}, и \code{empty}, и
порядок слов.

\begin{diybox}[Очередь в коде]
\textbf{Задание.} Запустить очередь и поймать переполнение.\\
\textbf{В Logra.} Вкладка «Код»: \code{fifo\_ptr.v} (или \code{fifo8.v}),
«Синтезировать», «Диаграммы»: раздел «События», импульсы \code{wr} по тактам
с \code{din}, затем \code{rd}.\\
\textbf{Ожидаемый результат.} После восьми записей \code{full = 1}, девятая
запись ничего не меняет; слова выходят в порядке записи; после последнего
чтения \code{empty = 1}.\\
\textbf{Как проверить.} \code{tb\_fifo\_ptr.v} печатает \code{fifo\_ptr: OK}.
\stars{2} Сделайте очередь на 4 слова по 4 бита, не меняя текста модуля.\\
\textbf{Образец ответа.} Параметры меняются при создании экземпляра:
\code{fifo\_ptr \#(.DW(4), .AW(2)) u (.clk(clk), .rst(rst), .wr(wr), .rd(rd),
.din(din), .dout(dout), .empty(empty), .full(full));}. Текст самого модуля
остаётся прежним; при \code{.AW(2)} признак \code{full} наступает после четырёх
записей (чтобы проверить это тестбенчем, замените в \code{tb\_fifo\_ptr.v}
число \code{8} на \code{4} в проверке и в модели).
\end{diybox}

\lookahead{в больших проектах очередь может быть огромной и хранить слова в
блочной памяти кристалла; бывает и \emph{двухтактная} очередь, где запись идёт
по одному тактовому сигналу, а чтение по другому~--- тогда указатели передают
между доменами в коде Грея (глава~\ref{ch:clocks}).}

\begin{warnbox}[Типичные ошибки]
\begin{itemize}
  \item \textbf{Различать пустую и полную по равенству указателей.} При
        $\text{wp}=\text{rp}$ непонятно, ноль слов или все. Нужен счётчик слов
        или лишний бит указателя.
  \item \textbf{Писать при \code{full} и читать при \code{empty}.} Слово
        теряется, либо читается мусор. Источник и приёмник должны смотреть на
        признаки.
  \item \textbf{Указатель не возвращается к нулю.} Если очередь на 8 слов
        использует 4-разрядный указатель без обнуления на границе, адрес
        выйдет за пределы памяти.
  \item \textbf{Забыть про сброс.} После включения указатели и счётчик
        должны быть нулями, иначе очередь «видит» несуществующие слова.
\end{itemize}
\end{warnbox}

\begin{diybox}[Очередь с признаками full и empty]
\textbf{Задание.} Проследить работу кольцевого буфера, найти моменты
«полна» и «пуста».\\
\textbf{В Logra.} Вкладка \emph{Примеры}, проект \code{fifo-four};
вкладки \emph{Диаграммы} и \emph{RTL}.\\
\stars{1} \emph{Повтори.} Дождитесь, пока очередь заполнится словами
1, 2, 3, 4, затем нажмите кнопку несколько раз; сверьте указатели с таблицей
трассы главы.\\
\stars{2} \emph{Измени.} Найдите на диаграмме такт, на котором \code{used}
достигает 4 и запись прекращается (так проявляется признак «полна»), и такт,
на котором \code{used} становится 0 («пуста»). Задача~5: что было бы без
защиты?\\
\stars{3} \emph{Проект.} На вкладке \emph{Код} увеличьте очередь до восьми
слов (трёхразрядные указатели, четырёхразрядный \code{used}, сравнение с 8) и
проверьте признаки полной и пустой очереди.\\
\textbf{Ожидаемый результат.} При полной очереди запись не меняет ни
указатель, ни счётчик; слова выходят в порядке записи.\\
\textbf{Как проверить.} Вкладка \emph{Диаграммы}: последовательность
\code{used} при заполнении и чтении совпадает с рис.~\ref{fig:fifo_flags};
вкладка \emph{RTL}: найдите память слов и два указателя.
\end{diybox}

\section{Вопросы и задачи}
\label{sec:tasks-fsmdesign}

Ответы собраны в приложении~\ref{app:answers}.

\begin{taskbox}[Задачи к главе]
\begin{enumerate}
  \item \stars{2}~Нарисуйте граф автомата Мура с тремя состояниями для светофора, у
        которого красный горит 4 такта, жёлтый 1, зелёный 3. Сколько триггеров
        нужно при двоичном кодировании? При унитарном?
  \item \stars{2}~Автомат по продаже получает одновременно монету $c5$ и $c10$ (ошибка
        датчика). Как поведёт себя описание из главы? Как изменить таблицу
        переходов, чтобы это учитывалось?
  \item \stars{2}~Изобразите таблицу переходов детектора последовательности \code{110}
        (с перекрытием). Автомат какого типа получится, Мура или Мили?
  \item \stars{1}~Очередь FIFO имеет 16 слов. Какой разрядности нужны указатели и
        счётчик слов?
  \item \stars{1}~В очередь на 8 слов подряд записали 11 слов. Сколько слов будет
        потеряно и почему?
  \item \stars{1}~Чем автомат Мили удобнее автомата Мура и чем опаснее?
  \item \stars{2}~\emph{Найдите ошибку.} Разработчик очереди на 8 слов определил
        $\text{full}=(\text{wp}=\text{rp})$ и $\text{empty}=(\text{wp}=\text{rp})$. В каких
        двух ситуациях признаки дадут неверный результат и как их исправить?
  \item \stars{2}~\emph{Найдите ошибку.} В таблице переходов автомата панели
        (раздел «Сквозной проект») проектировщик оставил пустой ячейку
        $d{=}1,c{=}1$ для состояния ОХРАНА, решив, что «этого не бывает». Опишите
        сценарий, в котором это случится, и скажите, к чему приведёт пустая
        ячейка. Какой переход вы запишете?
  \item \stars{3}~\emph{Мини-проект.} Соберите автомат режимов «Охранная панель»
        с вводом кода \code{1-0-1} (этап~9). Критерий приёмки: на диаграмме
        после \code{101} панель в состоянии ОХРАНА; $d=1$ в режиме ОХРАНА
        включает сирену со следующего такта; $d=1$ в режиме ОХРАНА\_ВЫКЛ её не
        включает; повторный ввод \code{101} в состоянии ТРЕВОГА её выключает.
  \item \stars{3}~\emph{Мини-проект.} Расширьте очередь на 4 слова выходами
        \code{full} и \code{empty} и сделайте так, чтобы запись при полной
        очереди была невозможна. Критерий приёмки: при пяти подряд записях
        на диаграмме \code{used}$=4$, указатель записи остановился, пятое слово не
        затёрло первое; при пяти подряд чтениях пятое чтение ничего не меняет.
\end{enumerate}
\end{taskbox}

\section{Что дальше}

Мы научились описывать поведение устройства автоматом и хранить данные в
очереди, то есть проектировать схему любой сложности и проверять её в Logra. Но
из чего такая схема сделана внутри микросхемы? В ПЛИС нет отдельных вентилей:
комбинационную логику заменяют таблицы, автомат и очередь собираются из
триггеров и блоков памяти, а провода между ними переключаются при загрузке
конфигурации. Об этом глава~\ref{ch:inside}: мы заглянем внутрь кристалла, а
затем разберём плату и маршрут от текста до прошивки.
